\section{Strength: where the force comes from and how a support structure takes it up}\label{sec:festigkeit}

The first structural calculation of the 1-GW plant (fusion surrogate with 79 samples, hex model and Elmer) gave stresses in the coil casing of
about \SI{2000}{MPa}, four times more than steel at \SI{4}{K} can bear in the long term. The user noted on this:
\emph{“Wall, coils and strength are one topic, they belong together. The plasma belongs to it too. The coil shape
determines the plasma, and the coil also has to hold, and that also affects the temperature.” (translated)} Later the user
added the package (everything between plasma and coil) and a 3D model as the goal, and when the numbers came:
\emph{“We have to concentrate on strength. 2000~MPa is a killer, even 500~MPa is a lot. Maybe we need more coils, or we
have to understand how the force arises.” (translated)} This section describes how, step by step, a design emerged
that keeps within the limits, under which conditions it holds, and why it works. The design study is complete and
evaluated; a separate variant report presents and compares all designs computed, this chapter shows the best ones and
those close to the limit.

\paragraph{Terms used in this section.}
\begin{description}\itemsep2pt
\item[Lorentz force] the force that a magnetic field exerts on a conductor carrying current. Per metre of conductor it
  is $\vec f = I\,\vec t\times\vec B$ (current $I$, direction of the conductor $\vec t$, field $\vec B$). In our coils
  \SI{22}{MA} flows; that gives up to \SI{140}{MN} per metre, the weight of \SI{14000}{t} on every metre.
\item[self-field] the field that a coil produces at its own conductor. It pushes the coil apart like an inflated tube.
\item[equivalent stress (von Mises)] a number that condenses the multiaxial stress state in the material into a value
  that can be compared with the strength from the tensile test.
\item[99-\% stress $\sigma_{99}$] (formerly called ``p99'') the stress that is not exceeded in 99~\% of the steel volume. Single peaks at edges of the
  coarse calculation mesh are thus not overrated.
\item[rope tension] A curved rope can carry a load pressing transversely on it by tension alone, if the load matches
  the curvature (like a ring under internal pressure). A straight piece cannot do this: it has to carry the load by
  \emph{bending}, and bending produces much higher stresses than tension.
\item[bearing, bucking cylinder] the places where the coils are held. In the tokamak each coil presses with its
  straight inner leg against a bucking cylinder in the centre of the machine.
\item[connector] steel blocks between neighbouring coil casings that pass forces from coil to coil.
\item[strain] the relative change in length. The superconducting tape (REBCO) tolerates only a few tenths of a percent
  of tension in the long term before it is damaged.
\item[package] all layers between plasma and coil: first wall, blanket (where tritium and heat are produced), shielding
  against neutrons and vacuum vessel.
\item[design study (DOE)] many designs whose control variables are spread over the whole range by plan (here Latin
  hypercube \cite{mckay1979}), so that one can see which quantity has which effect.
\end{description}

\paragraph{Division of labour.} The user set the direction and decided at every fork. The AI (Claude Opus) planned,
set expectations, wrote the coupling and evaluated the results. Subtasks went to smaller models: Claude Sonnet
researched the limit values (every source checked via Crossref, sample re-checked by Opus) and extended the coil
optimisation by a force target; Claude Haiku explored how a new coil set runs through the structural model. The
calculations were done with simsopt \cite{simsopt} (coils) and Elmer \cite{elmer} (strength).

\subsection*{Expectation beforehand}
Before each calculation, the event log recorded what was expected and how the result was to be interpreted:
\begin{enumerate}\itemsep1pt
\item A force target in the coil optimisation lowers the peak force by about 30~\% at double the field error.
\item Reducing force alone does not give the required factor~4; what is needed is a support structure that passes the
  forces on in a closed way.
\item The force per metre lies at 100--150~MN/m; the required rope tension varies along the coil by more than a factor
  of~3, because the curvature changes strongly. Bending, not tension, is the cause of the high stress.
\item The net force on each coil points to the centre of the machine.
\item More coils lower the force per metre roughly in inverse proportion to the number of coils; the steel needed for
  the same stress stays about the same.
\item The package is tight at the present spacing.
\end{enumerate}

\subsection*{Step 1: Force target in the coil optimisation}
The coil optimisation (stage~2) received an additional term that penalises high force per metre (simsopt
\texttt{LpCurveForce}). Without this term exactly the previous coil set results (deviation of the field error
\SI{0.01}{\%}).
\begin{center}\small
\begin{tabular}{lccc}
\toprule
Weight of force target & Field error $B\cdot n$ (mean) & Peak force & Mean force \\
\midrule
0 (previous) & 0.27~\% & 100~\% & 100~\% \\
0.01 & 0.30~\% & $-9$~\% & $-3$~\% \\
1 & 1.56~\% & $-25$~\% & $-6$~\% \\
\bottomrule
\end{tabular}
\end{center}
\textbf{Result:} The peaks can be capped, the mean force hardly. It is fixed by field strength and the size of the
machine. Expectation~1 was too optimistic, expectation~2 is confirmed: the coil shape alone does not solve the problem.
The force weight is only used further in the range 0 to 0.03.

\subsection*{Step 2: A coupled calculation chain}
So that plasma, coils, package, strength and temperature can be evaluated together, a chain was built
(\texttt{stellarator/kopplung.py}). A design runs through it in 1--3~minutes:
\begin{enumerate}\itemsep1pt
\item \textbf{Coils} (simsopt): number of coils, distance to the plasma, coil length, force target $\to$ field error as
  a measure of plasma quality.
\item \textbf{Calibration}: current such that $B_0 = \SI{6.5}{T}$ prevails on the magnetic axis. Winding pack square,
  edge length such that the current density stays as with 16 coils.
\item \textbf{Strength} (hex mesh, Elmer): casing with four wall thicknesses, potting, connectors, support; load
  $\vec f = I\vec t\times\vec B$ from all coils.
\item \textbf{Package}: smallest distance from the inside of the casing to the plasma edge, all round; from this the
  shield thickness, nuclear heating and neutron dose in the winding pack.
\item \textbf{Temperature}: field at the edge of the winding pack $\to$ utilisation of the superconductor and temperature
  margin.
\end{enumerate}
The limits and assumptions:
\begin{center}\small
\begin{tabular}{p{0.34\linewidth}p{0.18\linewidth}p{0.38\linewidth}}
\toprule
Quantity & Limit & Origin (status) \\
\midrule
Casing steel 1.4429/316LN, $\sigma_{99}$ & \SI{500}{MPa} & 2/3 of the yield strength at 4~K, ITER grade C3; best grade C1 allows 667~MPa (ITER slides, unverified) \\
Strain winding pack & 0.3~\% & REBCO irreversible from 0.45--0.72~\% \cite{barth2015}; ARC kept $<0.2$~\% \cite{sorbom2015} \\
Nuclear heating winding pack & \SI{2000}{W/m^3} & ARIES-CS (unverified); strict \SI{50}{W/m^3} for EU-DEMO \cite{cufar2019} \\
Neutron dose REBCO & $3\cdot10^{18}$~n/cm$^2$ & lower bound after ARC \cite{sorbom2015} \\
Utilisation of superconductor & 70~\% & assumption; $J_c(B,T)$ after \cite{senatore2024,pardo2026} \\
Field error $B\cdot n$ & 0.3~\% & previous coil set as benchmark \\
Package at least & \SI{1.0}{m} & ARIES-CS design \cite{elguebaly2005} \\
Neutron attenuation & e-fold per \SI{0.1}{m} & assumption; tungsten alone \SI{0.061}{m} \cite{windsor2021} \\
\bottomrule
\end{tabular}
\end{center}
While building the chain an old error came to light: the previous structural model had scaled the coils 3~\% larger than
the plasma and alpha calculation. The chain now uses the same scale everywhere. The first test point (previous coils)
confirmed the expectations with one exception: the field at the conductor came out at \SI{6.4}{T} instead of about
\SI{15}{T}. Cause: it had been measured on the centreline, where the self-field cancels. At the edge of the winding pack
it is \SI{18.7}{T}. The test point showed: package (\SI{1.77}{m} free), temperature (utilisation 34~\%) and field
accuracy fit. It fails only on strength ($\sigma_{99}$ = \SI{2129}{MPa}) and on strain in the winding pack (0.76~\%). Expectation~6
(package tight) was wrong: ARIES-CS manages with \SI{1.3}{m} from plasma to coil centre \cite{elguebaly2005}, we have
\SI{2.5}{m}.

\subsection*{Step 3: Where the force comes from}
For each coil the force per metre along the conductor was calculated: the field of the other coils by Biot--Savart, the
self-field with the regularised formula for a rectangular cross-section (simsopt). It is split into a part pointing
outwards (the coil widens) and a sideways part.

\bild{kraft_analyse_16_en.png}{1.0}{Force along the four basic coils (16 coils in total, the others are mirror images):
force per metre, required rope tension, curvature.}
\lesehilfe{Each column is one coil, horizontally the path along the coil (one coil is about \SI{45}{m} long). Top: black
is the total force per metre, blue the part that widens the coil, red the sideways part. The force is everywhere
\SI{50}{MN/m} or more and reaches \SI{140}{MN/m} in the middle of the path: this is the inner leg of the coil on the
inside of the ring, where the coils stand close together and the field is strongest. Middle: the rope tension with
which the coil could carry this load as pure tension. Where the curve is missing, the coil is almost straight; there
this would only work with infinite rope tension. Bottom: the curvature. Peaks are tight bends (radius about
\SI{2.5}{m}), valleys are almost straight pieces. The picture shows the core problem: the largest force sits exactly on
the straight piece.}

\textbf{Derivation: why straight pieces are so costly.} For a rope with rope tension $T$, curvature $\kappa$ and
transverse load $f_N$, equilibrium gives
\[
\frac{dT}{ds}\,\vec t + T\kappa\,\vec N + \vec f = 0 \quad\Rightarrow\quad T = \frac{f_N}{\kappa}.
\]
With $f_N \approx \SI{100}{MN/m}$ and $\kappa \approx \SI{0.25}{m^{-1}}$ one needs $T \approx \SI{400}{MN}$. With about
\SI{1.3}{m^2} of steel cross-section in the casing, this is already about \SI{300}{MPa}. This is the share that cannot be
avoided if field and size are fixed. On the inner leg the coil is almost straight over about $l = \SI{7}{m}$
($\kappa < \SI{0.08}{m^{-1}}$), the load there is $q \approx \SI{130}{MN/m}$. As a beam on two supports this gives
\[
M = \frac{q\,l^2}{8} \approx \SI{800}{MNm},\qquad
W = \frac{H^4 - h^4}{6H} \approx \SI{0.38}{m^3}\ (H = \SI{1.36}{m}\ \text{casing},\ h = \SI{0.755}{m}\ \text{winding pack}),
\qquad \sigma_B = \frac{M}{W} \approx \SI{2100}{MPa}.
\]
This back-of-the-envelope estimate matches the FEM value of \SI{2129}{MPa} remarkably well. The high stresses therefore
come from the bending of the straight inner leg.

\textbf{How tokamaks solve this.} A tokamak coil has the shape of a “D”: the arc is shaped so that it carries the load
as pure tension, and the straight leg on the inside lies flat against a bucking cylinder or against the neighbouring
coils (wedging). It is therefore not bent but compressed. Our model, in contrast, held each coil only at its outermost
point. The net force per coil, however, points inwards with up to \SI{1.1}{GN}, to the centre of the machine
(expectation~4 confirmed). The coil therefore hung like a beam supported on the wrong side.

\textbf{More coils?} The same analysis for 20 and 24 coils:
\begin{center}\small
\begin{tabular}{lcccc}
\toprule
Coils & Current per coil & Force per metre (mean) & Sum of all coils & Field error $B\cdot n$ \\
\midrule
16 & \SI{22.1}{MA} & \SI{85}{MN/m} & \SI{58}{GN} & 0.27~\% \\
20 & \SI{17.7}{MA} & \SI{65}{MN/m} & \SI{57}{GN} & 0.074~\% \\
24 & \SI{14.7}{MA} & \SI{54}{MN/m} & \SI{58}{GN} & 0.028~\% \\
\bottomrule
\end{tabular}
\end{center}
The force per metre falls, the sum stays (expectation~5 confirmed): with the same steel share the same stress follows.
The gain lies elsewhere: with more coils the field becomes ten times more accurate. This margin can be traded for
shorter, smoother coils.

\subsection*{Step 4: Support-structure trials with the same 16 coils}
Step~3 directly showed what has to change: support where the force points, and couple the coils into a composite that
passes transverse loads from coil to coil. For this, the structural model received a bearing at the innermost point of
each coil (face towards the centre of the machine, like a bucking cylinder), and optionally inside and outside.

\bild{festigkeit_tragwerk_en.png}{1.0}{Same 16 coils, different support structure: stress, deformation, strain and steel
mass for five variants (A--E), including the stress seen from above for A and E.}
\lesehilfe{Top: four bar charts for the five variants A to E (key in the line below). Red bars lie above the limit
(dashed), blue ones below it, grey quantities have no fixed limit. From A to E the stress in the steel falls from
\SI{2129}{MPa} to \SI{481}{MPa} and the deformation from \SI{399}{mm} to \SI{41}{mm}; the steel mass rises from
\num{15600} to \num{37000}~t in exchange. Bottom left: variant~A seen from above; red marks places above \SI{1200}{MPa},
mainly at the inner legs around the hole in the middle. Bottom right: variant~E; almost everything is blue (below
\SI{400}{MPa}); the coils have become thicker and, through connectors, a composite.}

\begin{center}\small
\begin{tabular}{llcccc}
\toprule
 & Variant & $\sigma_{99}$ steel & Deformation & Strain winding pack & Total steel \\
\midrule
A & bearing outside (previous), 6 connectors per pair & \SI{2129}{MPa} & \SI{399}{mm} & 0.76~\% & \SI{15.6}{kt} \\
B & bearing inside & \SI{1454}{MPa} & \SI{273}{mm} & 0.55~\% & \SI{15.6}{kt} \\
C & inside and outside & \SI{1372}{MPa} & \SI{226}{mm} & 0.54~\% & \SI{15.6}{kt} \\
D & + 12 connectors per pair, steel (200~GPa) & \SI{947}{MPa} & \SI{107}{mm} & 0.41~\% & \SI{23.1}{kt} \\
E & + thickest walls of the range & \SI{481}{MPa} & \SI{41}{mm} & 0.31~\% & \SI{37.0}{kt} \\
\bottomrule
\end{tabular}
\end{center}
\textbf{Why this works.} B: the net force goes into the bearing over a short path, instead of bending the whole coil
($-32$~\%, expectation 30--50~\% confirmed). D: more and stiffer connectors pass the transverse load of the straight leg
on to the neighbours; the free length $l$ of the beam becomes shorter, and $M \propto l^2$. E: a larger section modulus
$W$ lowers the bending stress further. With this, for the first time a design lies below \SI{500}{MPa}. The price is
mass. The problem has shifted: from “does not hold” to “too heavy”.

\subsection*{Step 5: Design study}
To find the most robust admissible design, 200 designs were computed (Latin hypercube). Eleven quantities were varied:
number of coils (16, 20, 24), coil--plasma distance, length per coil, force weight, four wall thicknesses, size of the
bearing areas, number (6--14 per pair) and stiffness (20--200~GPa) of the connectors. Bearings inside and outside are
fixed. During the study the user changed the goal: \emph{“The weight is not so important. It is a building, after
all.” (translated)} Since then the aim is not the lightest but the most robust design, i.e.\ the one with the largest
smallest reserve over all limits. One design failed in meshing, 199 were computed completely.

\bild{doe1_uebersicht_en.png}{1.0}{Design study: 99-\% stress in the case steel against six design variables, all 199 designs.}
\lesehilfe{Each circle is a computed design; the colour shows the number of coils (blue 16, orange 20, green 24),
filled circles meet all eight limits. Vertically the stress (logarithmic), dashed the limit of \SI{500}{MPa}. No single
variable orders the points; the clearest trend is the rise of stress with the coil--plasma distance (bottom centre).}

\paragraph{A clean-up was necessary.} Eleven designs did not fit into memory and were computed with a coarser mesh.
Only the mesh study revealed that their connectors became one third longer, because a connector always spans four
stations. These designs were therefore too favourable. Recomputed with the standard mesh, two of the originally nine
admissible designs drop out (0063: \SI{369}{\to}\SI{595}{MPa}, 0136: \SI{479}{\to}\SI{577}{MPa}).
\textbf{Seven admissible designs remain}, all with 20 or 24 coils.

\paragraph{List~1: the best and those close to the limit.} The table shows the admissible designs ranked by their
smallest reserve and the twelve that fail most narrowly. All other designs and a catalogue of all geometrically
representable ones are in the variant report.
\begin{table}[H]\centering\scriptsize\begin{tabular}{@{}lrrrrrrrl@{}}\toprule
No. & coils & $\sigma_{99}$ & max. & strain & $B\cdot n$ & clearance & steel & smallest reserve (bottleneck) \\
 & & MPa & MPa & \% & ‰ & m & kt &  \\\midrule
\multicolumn{9}{@{}l}{\textit{admissible}} \\
0061 & 24 & \num{351} & \num{725} & \num{0.17} & \num{1.49} & \num{1.62} & \num{19} & \num{30}\,\% (steel σ99) \\
0169$^\ast$ & 24 & \num{308} & \num{717} & \num{0.15} & \num{2.13} & \num{1.39} & \num{28} & \num{29}\,\% (package space) \\
0093 & 24 & \num{283} & \num{1190} & \num{0.16} & \num{1.45} & \num{1.37} & \num{28} & \num{27}\,\% (package space) \\
0008 & 24 & \num{305} & \num{1695} & \num{0.17} & \num{2.20} & \num{1.48} & \num{26} & \num{27}\,\% (field error B·n) \\
0014 & 20 & \num{377} & \num{1419} & \num{0.21} & \num{2.21} & \num{1.65} & \num{28} & \num{25}\,\% (steel σ99) \\
0017 & 24 & \num{389} & \num{1721} & \num{0.20} & \num{1.66} & \num{1.30} & \num{23} & \num{20}\,\% (package space) \\
0157 & 20 & \num{492} & \num{1658} & \num{0.22} & \num{2.37} & \num{1.35} & \num{18} & \num{2}\,\% (steel σ99) \\
\multicolumn{9}{@{}l}{\textit{narrowly missed}} \\
0173 & 24 & \num{533} & \num{1827} & \num{0.24} & \num{1.85} & \num{1.28} & \num{27} & \num{-7}\,\% (steel σ99) \\
0103 & 20 & \num{440} & \num{738} & \num{0.19} & \num{2.66} & \num{1.25} & \num{21} & \num{-8}\,\% (neutron fluence) \\
0044 & 20 & \num{540} & \num{1583} & \num{0.23} & \num{2.76} & \num{1.35} & \num{26} & \num{-8}\,\% (steel σ99) \\
0004 & 20 & \num{558} & \num{1085} & \num{0.26} & \num{1.81} & \num{1.54} & \num{19} & \num{-12}\,\% (steel σ99) \\
0136$^\ast$ & 20 & \num{577} & \num{2238} & \num{0.30} & \num{1.62} & \num{1.84} & \num{26} & \num{-15}\,\% (steel σ99) \\
0022 & 20 & \num{600} & \num{2183} & \num{0.26} & \num{0.70} & \num{1.57} & \num{24} & \num{-20}\,\% (steel σ99) \\
0137 & 24 & \num{612} & \num{2562} & \num{0.29} & \num{0.94} & \num{1.57} & \num{16} & \num{-22}\,\% (steel σ99) \\
0063$^\ast$ & 20 & \num{595} & \num{4265} & \num{0.37} & \num{0.75} & \num{1.53} & \num{41} & \num{-24}\,\% (winding-pack strain) \\
0042 & 20 & \num{632} & \num{3962} & \num{0.34} & \num{2.66} & \num{1.61} & \num{21} & \num{-26}\,\% (steel σ99) \\
0043 & 24 & \num{653} & \num{1208} & \num{0.29} & \num{1.79} & \num{1.30} & \num{16} & \num{-31}\,\% (steel σ99) \\
0140 & 20 & \num{656} & \num{4169} & \num{0.34} & \num{1.46} & \num{1.76} & \num{22} & \num{-31}\,\% (steel σ99) \\
0085 & 20 & \num{608} & \num{2067} & \num{0.28} & \num{2.00} & \num{1.23} & \num{30} & \num{-31}\,\% (neutron fluence) \\
\bottomrule\end{tabular}

\caption{Ranking (cleaned). $^\ast$ = recomputed with the standard mesh.}\end{table}

\bild{doe1_reserven_en.png}{0.95}{Reserves of the admissible designs per limit.}
\lesehilfe{Rows are designs, columns limits. The number is the distance to the limit in per cent, darker means more
reserve; bold is the bottleneck. For the best designs, steel stress, package space and field error alternate as the
bottleneck; superconductor and nuclear heating have large reserves everywhere.}

\subsection*{Step 6: What determines the stress}
For six quantities a surrogate model (Gaussian process \cite{rasmussen2006}) was trained and checked by five-fold
cross-validation. Package space and nuclear heating are predicted almost perfectly ($R^2=0.98$), steel mass and field
error well (0.95 and 0.90), the steel stress poorly ($R^2=0.46$). Its most important variable is the coil--plasma
distance, then the coil length. Number and stiffness of the connectors matter less than expected. Interpretation: the
stress depends on the coil shape -- coils further away must be more strongly twisted and carry more current -- and on
the random position of the connectors, which sets the local peaks.
\bild{doe1_ersatz_en.png}{1.0}{Surrogate models: prediction against computation (top) and influence of the variables (bottom).}
\lesehilfe{Top: each point is a design, horizontally computed, vertically predicted by a model that had not seen it;
on the diagonal the prediction would be perfect. Bottom: how strongly each variable moves each quantity. The field
error depends on the number of coils, package and heating on the distance, the mass on connector count and wall
thickness.}

\paragraph{Putting it to the test.} The surrogate proposed four new designs that maximise the smallest reserve minus
one standard deviation. When computed, three were admissible and one failed clearly ($\sigma_{99}$ \SI{851}{MPa}).
The best proposal V3 reached a robustness of 0.35 on the standard mesh, more than 0169, but lost on the fine mesh
($\sigma_{99}$ \SI{272}{} $\to$ \SI{394}{MPa}) and dropped to 0.21. The surrogate is therefore useful for
pre-selection, but stress still has to be computed directly; 0169 remains the most robust design.

\subsection*{Step 7: Mesh study -- how certain is ``admissible''?}
The five best designs were recomputed with the same coils and a finer mesh: 1.5 times as many stations along the
coil (same connector length) or three instead of two cells across the wall.
\begin{table}[H]\centering\small\begin{tabular}{@{}lrrrrrrrr@{}}\toprule
Design & \multicolumn{2}{c}{doe1} & \multicolumn{2}{c}{standard} & \multicolumn{2}{c}{1.5× stations} & \multicolumn{2}{c}{wall 3 cells} \\
  & $\sigma_{99}$ & max. & $\sigma_{99}$ & max. & $\sigma_{99}$ & max. & $\sigma_{99}$ & max. \\\midrule
0008 & \num{305} & \num{1695} & -- & -- & \num{407} & \num{3250} & \num{306} & \num{1931} \\
0014 & \num{377} & \num{1419} & -- & -- & \num{549} & \num{3814} & \num{375} & \num{1639} \\
0018 & \num{811} & \num{2650} & \num{996} & \num{4566} & -- & -- & -- & -- \\
0055 & \num{366} & \num{753} & \num{413} & \num{942} & -- & -- & -- & -- \\
0060 & \num{662} & \num{1313} & \num{704} & \num{1499} & -- & -- & -- & -- \\
0061 & \num{351} & \num{725} & -- & -- & \num{556} & \num{6553} & \num{346} & \num{861} \\
0063 & \num{369} & \num{2016} & \num{595} & \num{4265} & -- & -- & -- & -- \\
0084 & \num{395} & \num{911} & \num{462} & \num{1018} & -- & -- & -- & -- \\
0093 & \num{283} & \num{1190} & -- & -- & \num{410} & \num{2343} & \num{280} & \num{1432} \\
0121 & \num{264} & \num{520} & \num{286} & \num{749} & -- & -- & -- & -- \\
0136 & \num{479} & \num{995} & \num{577} & \num{2238} & -- & -- & -- & -- \\
0145 & \num{362} & \num{1838} & \num{426} & \num{2490} & -- & -- & -- & -- \\
0149 & \num{837} & \num{1541} & \num{636} & \num{1423} & -- & -- & -- & -- \\
0169 & \num{219} & \num{421} & \num{308} & \num{717} & \num{319} & \num{1262} & \num{308} & \num{794} \\
0176 & \num{245} & \num{615} & \num{434} & \num{2674} & -- & -- & -- & -- \\
\bottomrule\end{tabular}
\caption{Mesh study: $\sigma_{99}$ and
largest element value in the steel (MPa).}\end{table}
\bild{doe1_netz_en.png}{0.95}{Mesh study of the best designs.}
\lesehilfe{Left $\sigma_{99}$, right the largest element value, finer from left to right. A flat line means the result
does not depend on the mesh. Finer walls change almost nothing; more stations along the coil raise $\sigma_{99}$ by
34--58~\% for four designs. Only 0169 stays nearly the same (\SI{308}{\to}\SI{319}{MPa}). The maxima grow with every
refinement: the sign of a sharp edge.}
Of the 200 most highly stressed steel elements, 77--100~\% lie less than \SI{1}{m} from a connector joint. The
connectors are sharp-edged steel blocks; at their ends along the coil the stiffness jumps, as at the end of a lap
joint. \textbf{Only design 0169 is robust against the mesh} (24 coils, 13 stiff connectors per pair, \SI{28}{kt}
steel); it is therefore the best design of this study. The others are ``admissible'' only on the standard mesh.

\subsection*{Step 8: Which steel, which limit?}\label{sec:stahl}
The user asked: \emph{“I want to know exactly which steel you are using. I don't know p99 steel.” (translated)} The case
steel is \textbf{1.4429 / X2CrNiMoN17-13-3} (AISI 316LN): austenitic, about 17~\% chromium, 13~\% nickel, 2.5~\%
molybdenum, low carbon, strengthened with nitrogen; ITER builds its coil cases from it. ``p99'' was not a material but
the quantity $\sigma_{99}$ (see terms).

\bild{stahl_kennwerte_en.png}{1.0}{1.4429 / 316LN at room temperature and at \SI{4}{K}.}
\lesehilfe{Per group blue is the yield strength, green the tensile strength, yellow two thirds of the yield strength.
On the left the handbook value for room temperature (\SI{295}{MPa} \cite{schweizerfn}), next to it measurements: at
\SI{4.2}{K} 316LN reaches \SI{1070}{MPa} yield and \SI{1543}{MPa} tensile strength \cite{fedorikova2023}, with high
nitrogen above \SI{1000}{MPa} even in thick plates \cite{sakurai2024}. On the right the guaranteed ITER minimum values.}
\textbf{Derivation:} The ITER magnet criteria allow $S_m=\tfrac23R_{p0.2}$ at operating temperature. With the weakest
ITER grade C3 ($R_{p0.2}\ge\SI{750}{MPa}$ at \SI{4}{K}) this gives \SI{500}{MPa}; the ASME form
$\min(\tfrac23R_{p0.2},\tfrac13R_m)$ with the measured 4-K values gives \SI{514}{MPa}. At room temperature only about
\SI{200}{MPa} would be allowed; this does not matter, because the magnetic force acts only in cold operation.

\paragraph{Stress categories and maxima.} In addition to $\sigma_{99}$, the categories according to ITER/ASME were
determined: for each section across the wall the mean (membrane stress $P_m$, allowable $S_m$) and the linear part
(bending $P_b$, allowable $P_m+P_b\le1.5\,S_m$).
\bild{doe1_kategorien_en.png}{1.0}{Stress categories and maxima.}
\lesehilfe{Two bars per design: membrane stress $P_m$ (blue) and membrane+bending (orange); dashed the limit
\SI{500}{MPa} for $P_m$, dotted \SI{750}{MPa} for $P_m+P_b$. Left the maximum over the whole steel; the diamond marks
the largest single element: maximum and peak lie together, at the end of a connector -- a local discontinuity that
mostly does not even meet the limit $1.5\,S_m$. Right the value met by 99~\% of all wall sections: for the best
designs and the proposals V3 and V1 both bars stay below their limits, so the walls themselves can carry the load.}
Result: in 99~\% of the wall sections $P_m$ stays below \SI{400}{MPa} for the best designs; the walls can carry the
load. The connector joints cannot in the model; there the design needs a shape (fillet, tapered thickness, longer
contact area).

\subsection*{Step 9: Weakening selectively -- first attempt}
The user's idea (\emph{“It can also help to weaken certain areas so that the load is distributed more evenly.”
(translated)}) was first tested in its simplest form: all connectors of 0169 softer.
\begin{table}[H]\centering\scriptsize\begin{tabular}{@{}lrrrrrrrl@{}}\toprule
No. & E connectors & stations & $\sigma_{99}$ & max. & deformation & strain & $B\cdot n$ & all limits \\
 & & & MPa & MPa & mm & \% & ‰ & \\\midrule
0105 & 50 GPa & 64 & \num{358} & \num{743} & \num{27} & \num{0.17} & \num{2.13} & yes \\
0110 & 100 GPa & 64 & \num{331} & \num{734} & \num{22} & \num{0.16} & \num{2.13} & yes \\
0114 & 140 GPa & 64 & \num{318} & \num{725} & \num{21} & \num{0.15} & \num{2.13} & yes \\
0205 & 50 GPa & 96 & \num{372} & \num{1541} & \num{29} & \num{0.19} & \num{2.13} & yes \\
0210 & 100 GPa & 96 & \num{343} & \num{1398} & \num{25} & \num{0.18} & \num{2.13} & yes \\
0214 & 140 GPa & 96 & \num{330} & \num{1323} & \num{23} & \num{0.17} & \num{2.13} & yes \\
\bottomrule\end{tabular}
\caption{0169 with softer
connectors; numbers 1xx standard mesh, 2xx fine mesh (original: \SI{183}{GPa}).}\end{table}
It was expected (F-139) that the peaks at the connector ends would drop and $\sigma_{99}$ would fall by 5--15~\%. The
opposite happened: with \SI{50}{GPa} instead of \SI{183}{GPa}, $\sigma_{99}$ in the case rises from \SI{308}{MPa} to
\SI{358}{MPa} ($+16$~\%) and the deformation from \SI{19}{mm} to \SI{27}{mm}; only the connectors themselves are
relieved (\SI{192}{\to}\SI{127}{MPa}). The reason: here the connectors are the relief path. Making them soft overall
sends the load back into bending of the case walls. The idea remains right but must be applied \emph{locally}: soft
ends (where the peak sits) and a stiff core that keeps carrying the load. This needs connectors with graded stiffness in
the model -- the next step.


\subsection*{Step 10: The 3D model, entirely in PicoGK}
The user asked: \emph{“Can you please try to build a 3D model entirely in PicoGK?” (translated)} PicoGK \cite{picogk}
describes bodies with voxels and distance fields; union, subtraction and offset always succeed, even for the twisted
cases where Boolean operations in FreeCAD often fail. Python passes only geometric data (cross-sections per station,
connector surfaces, plasma boundary, bearing locations). In PicoGK the winding packs, cases (outer contour minus
potting) and potting (minus winding pack) are built as Boolean differences, the package by offsetting the plasma
(first wall and blanket to \SI{0.55}{m}, shield to \SI{0.88}{m}, vacuum vessel to \SI{1.0}{m}), the bearings as rods
to an inner support cylinder and an outer ring. Three test bodies show overlaps.
\bild{modell3d_0169_schnitt_en.png}{0.95}{Design 0169 as PicoGK model in a stepped cut.}
\lesehilfe{Each layer is cut open at a different angle. From inside: plasma (pink), first wall and blanket (yellow),
shield (green), vacuum vessel (grey). Around it the 24 coil cases (blue) with their winding pack (orange, visible in the
cut), connectors in between (dark blue). Dark grey the schematic bearings: rods from the coils to a support cylinder in
the centre and a ring outside.}
Winding pack and case agree with the computational model to 0.5--1~\%, the connectors mostly to 0--3~\% (expected
$\pm3$~\%). For 0169, 13~\% of the connector volume is missing: neighbouring connectors partly interpenetrate there,
and the computational model counts the shared space twice. Connector mass and stiffness of 0169 are therefore
somewhat too high in the computational model; according to the study with softer connectors this changes
$\sigma_{99}$ by a few percent, but the connectors should lie without overlap in a recomputation. Only the 3D model
revealed this error. PicoGK finds no overlap of
package and coils and none between the cases -- an independent confirmation of the clearance and penetration checks.
Only the \SI{2}{cm} thin potting cannot be represented with \SI{40}{mm} voxels. One model takes 2--5 minutes. Models
of all designs of list~1 and of all geometrically representable designs are in the variant report.

\subsection*{Comparison with the expectation}
\begin{center}\small
\begin{tabular}{p{0.45\linewidth}p{0.45\linewidth}}
\toprule
Expectation & Result \\
\midrule
1. Force target: $-30$~\% at double field error & more expensive: $-9$~\% at $+12$~\%, $-25$~\% at $\times5.7$ \\
2. Force reduction alone is not enough & confirmed \\
3. 100--150~MN/m, rope force varies $>\times3$, bending is the cause & confirmed; hand calculation matches FEM \\
4. Net force towards the plant centre & confirmed, up to \SI{1.1}{GN} per coil \\
5. More coils: same sum, same stress & confirmed; the gain is field accuracy \\
6. Package tight & refuted: 1.3--\SI{1.8}{m} free, ARIES-CS needs about \SI{1}{m} \\
Bearing inside: $-30$ to $-50$~\% & confirmed: $-32$~\% \\
DOE: 10--25~\% below \SI{500}{MPa} & confirmed: 11~\% \\
DOE: strain second bottleneck ($<10$~\% pass) & refuted: 24~\% pass \\
DOE: connectors and side wall set the stress & refuted: coil--plasma distance and coil length \\
Surrogate for stress $R^2\ge0.8$ & refuted: 0.46 (package, heating 0.98) \\
Finer mesh: $+5$ to $+25$~\% & refuted: $+34$ to $+58$~\% along the coil; only 0169 $+4$~\% \\
Softer connectors: $\sigma_{99}$ $-5$ to $-15$~\% & refuted: rises (see step~9) \\
PicoGK: volume $\pm3$~\%, no overlap & confirmed \\
\bottomrule
\end{tabular}
\end{center}

\subsection*{Limitations}
\begin{itemize}\itemsep1pt
\item \textbf{Rigid bearings.} The bearing areas are clamped as if the support cylinder were infinitely stiff. This is
  optimistic; the support cylinder itself must be modelled.
\item \textbf{Sharp-edged connector joints.} Mesh-dependent peaks arise there; the assessment applies to the walls, not
  to the joints.
\item \textbf{Primary load only.} Linear elastic, magnetic force only. Cool-down, dead weight, fatigue and fracture
  mechanics (often governing at \SI{4}{K}) are missing.
\item \textbf{Package and temperature coarse.} Neutron attenuation and tape fraction are assumptions.
\item \textbf{Plasma only via the field error.} Alpha losses for 0169 are still to be computed in full.
\end{itemize}

\subsection*{Next steps}
Shape the connector joints and make them compliant locally (not globally): soft ends, stiff core, iteration by stress
ratio; elastic support cylinder instead of rigid bearings; alpha orbits in the coil field of 0169; fatigue and fracture
mechanics at \SI{4}{K}.
